\par\Needspace{6\baselineskip}\medskip
\begin{center}\hypertarget{native-189-article141}{}\pdfbookmark[2]{141.}{article-141}141.\end{center}
Aus $P=\dfrac{n''}{n}$ und Gleichung III des vorigen Artikels folgt
\[
(n+n'')\frac{P+a}{P+1}=bn'\frac{\sin(z-\sigma)}{\sin z}\mathpunct{\raisebox{.45ex}{;}}
\]
Daraus und aus $Q=2\left(\dfrac{n+n''}{n'}-1\right)r'^3$ sowie $r'=\dfrac{R'\sin\delta'}{\sin z}$ ergibt sich
\[
\sin z+\frac{Q\sin z^4}{2R'^3\sin\delta'^3}=b\frac{P+1}{P+a}\sin(z-\sigma)
\]
oder
\[
\begin{aligned}
\frac{Q\sin z^4}{2R'^3\sin\delta'^3}&=\left(b\frac{P+1}{P+a}-\cos\sigma\right)\sin(z-\sigma)\\[.5ex]
&\quad-\sin\sigma\cos(z-\sigma).
\end{aligned}
\]
Setzen wir also der Kürze halber
\[
\tag*{14)}\frac{1}{2R'^3\sin\delta'^3\sin\sigma}=c
\]
und führen einen Hilfswinkel $\omega$ so ein, dass
\[
\operatorname{tang}\omega=\frac{\sin\sigma}{b\dfrac{P+1}{P+a}-\cos\sigma}\mathpunct{\raisebox{.45ex}{,}}
\]
so ergibt sich die Gleichung
\[
\tag*{IV.}cQ\sin\omega\sin z^4=\sin(z-\omega-\sigma),
\]
\SourceLeaf{190}aus der die Unbekannte $z$ zu bestimmen ist. Damit sich der Winkel $\omega$ bequemer berechnen lässt, schreibt man die vorige Formel für $\operatorname{tang}\omega$ zweckmäßig in der Form
\[
\operatorname{tang}\omega=\frac{(P+a)\operatorname{tang}\sigma}{P\left(\dfrac{b}{\cos\sigma}-1\right)+\left(\dfrac{b}{\cos\sigma}-a\right)}\mathpunct{\cdot}
\]
Setzt man daher
\[
\tag*{15)}\frac{\dfrac{b}{\cos\sigma}-a}{\dfrac{b}{\cos\sigma}-1}=d
\]
\[
\tag*{16)}\frac{\operatorname{tang}\sigma}{\dfrac{b}{\cos\sigma}-1}=e
\]
so erhält man zur Bestimmung von $\omega$ die folgende sehr einfache Formel
\[
\operatorname{tang}\omega=\frac{e(P+a)}{P+d}\mathpunct{\cdot}
\]
Die Berechnung der Größen $a,b,c,d,e$ nach den Formeln 11–16 hängt allein von den gegebenen Größen ab und bildet die vierte Aufgabe. Von den Größen $b,c,e$ werden nur die Logarithmen benötigt.

In einem Sonderfall müssen diese Vorschriften etwas abgeändert werden. Fällt nämlich der Großkreis $BB''$ mit $A''B''$ zusammen und damit die Punkte $B,B^\ast$ der Reihe nach mit $D',D$, so würden die Größen $a,b$ unendlich. Setzt man in diesem Fall
\[
\frac{R\sin\delta\sin(A'D''-\delta'+\sigma)}{R'\sin\delta'\sin(AD''-\delta)}=\pi,
\]
so erhält man anstelle von Gleichung III die Gleichung $0=\pi n-\dfrac{n'\sin(z-\sigma)}{\sin z}\mathpunct{\raisebox{.45ex}{,}}$ und indem man
\[
\operatorname{tang}\omega=\frac{\pi\sin\sigma}{P+(1-\pi\cos\sigma)}
\]
setzt, erhält man dieselbe Gleichung IV.

Ähnlich verhält es sich im Sonderfall $\sigma=0$. Dann wird $c$ unendlich und es gilt $\omega=0$, sodass der Faktor $c\sin\omega$ in Gleichung IV unbestimmt erscheint. Er ist jedoch bestimmt und hat den Wert
\[
=\frac{P+a}{2R'^3\sin\delta'^3(b-1)(P+d)}\mathpunct{\raisebox{.45ex}{,}}
\]
wie man leicht sieht. Damit gilt in diesem Fall
\[
\sin z=R'\sin\delta'\;\BareRoot[3]{\frac{2(b-1)(P+d)}{Q(P+a)}}\mathpunct{\cdot}
\]
